\documentclass[border=4pt]{standalone}
\usepackage{amsmath,amssymb,mathtools,xcolor,varwidth}
\definecolor{lineblue}{HTML}{1E6FE0}
\definecolor{curvegreen}{HTML}{00A35A}
\definecolor{ordorange}{HTML}{E8770A}
\begin{document}
\begin{varwidth}{28em}
\large\bfseries
A right line \textcolor{lineblue}{$AE$} and a \textcolor{curvegreen}{curve $ABC$} cut each other at the      given point $A$. From the curve points $B$ and $C$ we let fall the \textcolor{ordorange}{ordinates      $BD$ and $CE$}, drawn at a fixed angle to the base, so that two triangles $\triangle ABD$      and $\triangle ACE$ are formed beneath the curve. As $B$ and $C$ slide down toward $A$      both triangles dwindle to nothing, yet we may still ask the limiting value of their ratio.      Newton's claim is that this ultimate ratio is the \textcolor{ordorange}{duplicate ratio of the sides},      namely $AD^2 : AE^2$.
\end{varwidth}
\end{document}
